% !TEX program = pdflatex
\documentclass[11pt,a4paper]{article}

% =========================================================================
% PLAN Y REGISTRO DE RIESGOS
% Proyecto: CasaScore (AlquilaPuno)
% Gestión de Proyectos TICS (EST315) -- Semestre 2026-II -- Team Sobrevivientes
% =========================================================================

% --- PAQUETES DE CONFIGURACIÓN ---
\usepackage[utf8]{inputenc}
\usepackage[T1]{fontenc}
\usepackage[spanish,es-tabla]{babel}
% Compatibilidad con babel-spanish reciente (flechas de TikZ y signo de porcentaje)
\AtBeginDocument{\shorthandoff{<>}\spanishplainpercent}
\usepackage{geometry}
\geometry{
    top=3.6cm,
    bottom=2.3cm,
    left=2.2cm,
    right=2.2cm,
    % PARTE MODIFICABLE: headheight debe ser mayor que la altura del logo más alto (2.56cm = 73pt).
    headheight=80pt,
    headsep=15pt,
    footskip=20pt
}

\usepackage{graphicx}
\usepackage{fancyhdr}
\usepackage[table]{xcolor}
\usepackage{titlesec}
\usepackage{booktabs}
\usepackage{tabularx}
\usepackage{array}
\usepackage{amsmath}
\usepackage{enumitem}
\usepackage{microtype}
\usepackage{amssymb}
% --- RECURSOS VISUALES: GRÁFICOS Y CÓDIGO DE COLORES ---
\usepackage{tikz}
\usetikzlibrary{positioning,calc}
\definecolor{GrisFila}{RGB}{246, 248, 251}
\definecolor{CuadCerca}{RGB}{150, 30, 30}
\definecolor{CuadSatisf}{RGB}{180, 130, 20}
\definecolor{CuadInfo}{RGB}{10, 45, 90}
\definecolor{CuadMonit}{RGB}{110, 115, 125}
\newcommand{\etiqueta}[2]{{\setlength{\fboxsep}{2pt}\colorbox{#1}{\textcolor{white}{\scriptsize\bfseries #2}}}}
\newcommand{\cerca}{\etiqueta{CuadCerca}{Gestionar de cerca}}
\newcommand{\satisf}{\etiqueta{CuadSatisf}{Mantener satisfecho}}
\newcommand{\informado}{\etiqueta{CuadInfo}{Mantener informado}}
\newcommand{\monit}{\etiqueta{CuadMonit}{Monitorear}}
\usepackage{needspace}
\usepackage{longtable}
\definecolor{Rama1}{RGB}{10, 45, 90}
\definecolor{Rama2}{RGB}{40, 95, 150}
\definecolor{Rama3}{RGB}{180, 130, 20}
\definecolor{Rama4}{RGB}{30, 110, 60}
\definecolor{Rama5}{RGB}{110, 60, 140}
\definecolor{RaciR}{RGB}{150, 30, 30}
\definecolor{RaciA}{RGB}{10, 45, 90}
\definecolor{RaciC}{RGB}{180, 130, 20}
\definecolor{RaciI}{RGB}{140, 145, 155}
\definecolor{NivAlto}{RGB}{30, 110, 60}
\definecolor{NivMedio}{RGB}{180, 130, 20}
\definecolor{NivBajo}{RGB}{150, 30, 30}
\usetikzlibrary{patterns}
\definecolor{SevBaja}{RGB}{30, 110, 60}
\definecolor{SevMedia}{RGB}{180, 130, 20}
\definecolor{SevAlta}{RGB}{150, 30, 30}
\usetikzlibrary{arrows.meta,shapes.geometric,fit,backgrounds}
\usepackage{listings}
\lstset{basicstyle=\ttfamily\scriptsize, frame=single, rulecolor=\color{AzulUNA!40}, backgroundcolor=\color{GrisFila}, columns=fullflexible, keepspaces=true, breaklines=true, xleftmargin=4pt, xrightmargin=4pt, literate={á}{{\'a}}1 {é}{{\'e}}1 {í}{{\'i}}1 {ó}{{\'o}}1 {ú}{{\'u}}1 {ñ}{{\~n}}1 {Á}{{\'A}}1}

\usepackage{hyperref}

% Logos: funciona a 2 o 3 niveles de profundidad dentro de docs/latex/
\graphicspath{{../../assets/img/}{../../../assets/img/}}

% --- PALETA INSTITUCIONAL FIEI / UNA PUNO ---
\definecolor{AzulUNA}{RGB}{10, 45, 90}
\definecolor{DoradoFIEI}{RGB}{180, 130, 20}
\definecolor{GrisOscuro}{RGB}{50, 50, 50}
\definecolor{AzulClaro}{RGB}{228, 236, 246}

% --- CONFIGURACIÓN DE LISTAS ---
\setlist{topsep=2pt, itemsep=2pt, parsep=1pt}

% --- COLUMNAS DE TABLA ---
\newcolumntype{L}{>{\raggedright\arraybackslash}X}
\newcolumntype{P}[1]{>{\raggedright\arraybackslash}p{#1}}
\newcolumntype{C}[1]{>{\centering\arraybackslash}p{#1}}
\newcolumntype{Q}[1]{>{\raggedleft\arraybackslash}p{#1}}
\newcolumntype{R}{>{\raggedleft\arraybackslash}X}
\renewcommand{\arraystretch}{1.25}

% --- FORMATO DE ENLACES ---
\hypersetup{
    colorlinks=true,
    linkcolor=AzulUNA,
    urlcolor=AzulUNA,
    citecolor=AzulUNA,
    pdftitle={Plan y Registro de Riesgos -- CasaScore (AlquilaPuno)},
    pdfauthor={Team Sobrevivientes}
}

% --- FORMATO DE TÍTULOS ---
\titleformat{\section}{\color{AzulUNA}\normalfont\large\bfseries}{\thesection.}{0.7em}{}[\color{AzulUNA}\titlerule]
\titleformat{\subsection}{\color{AzulUNA}\normalfont\normalsize\bfseries}{\thesubsection.}{0.5em}{}
\titlespacing*{\section}{0pt}{9pt}{4pt}
\titlespacing*{\subsection}{0pt}{6pt}{2pt}

% --- ENCABEZADO Y PIE DE PÁGINA (fancyhdr) ---
\pagestyle{fancy}
\fancyhf{}

% ===== TAMAÑO DE LOS LOGOS (PARTE MODIFICABLE) =====
% logo_fiei.png  : height=1.45cm (lado izquierdo)
% logo_equipo.png: height=2.56cm (lado derecho; 25 % más grande porque la imagen tiene más borde blanco)
% Para cambiar el tamaño, modifique solo el valor height=... de cada \includegraphics.
% Si el logo más alto supera 2.6cm, aumente también headheight y top en la configuración de página.
\fancyhead[L]{%
    \includegraphics[height=1.45cm,keepaspectratio]{logo_fiei.png}%
}
\fancyhead[C]{%
    \scriptsize \textbf{\textcolor{AzulUNA}{UNIVERSIDAD NACIONAL DEL ALTIPLANO DE PUNO}}\\[2pt]
    \tiny \textcolor{GrisOscuro}{FACULTAD DE INGENIERÍA ESTADÍSTICA E INFORMÁTICA}\\[2pt]
    \scriptsize \textbf{\textcolor{AzulUNA}{Gestión de Proyectos TICS (EST315) \textbar\ Team Sobrevivientes}}\\[2pt]
    \footnotesize \textbf{\textcolor{DoradoFIEI}{PLAN Y REGISTRO DE RIESGOS}}%
}
\fancyhead[R]{%
    \includegraphics[height=2.56cm,keepaspectratio]{logo_equipo.png}%
}

\renewcommand{\headrulewidth}{0pt}
\renewcommand{\headrule}{}

\fancyfoot[L]{\scriptsize \textcolor{GrisOscuro}{CasaScore (AlquilaPuno)}}
\fancyfoot[C]{\small \thepage}
\fancyfoot[R]{\scriptsize \textcolor{GrisOscuro}{Puno, Perú -- Semestre 2026-II}}
\renewcommand{\footrulewidth}{0.4pt}

\begin{document}

\begin{center}
    {\Large \textbf{\textcolor{AzulUNA}{PLAN Y REGISTRO DE RIESGOS}}}\\[3pt]
    {\normalsize \textbf{CasaScore (AlquilaPuno)}}\\[4pt]
    \footnotesize \textbf{Fecha:} martes 22/09/2026 (semana 5) \quad \textbar\quad \textbf{Versión:} 1.0 \quad \textbar\quad \textbf{Fase:} Planificación
\end{center}

\vspace{0.1cm}

\section{Datos Generales}
\noindent\begin{tabularx}{\textwidth}{@{}P{4.2cm} L@{}}
\toprule
\textbf{Proyecto:} & CasaScore (aplicación AlquilaPuno) \\
\textbf{Responsable:} & Viviana Mary Quisbert Quispe -- Gestora de Comunicación, Riesgos y Documentación \\
\textbf{Proceso PM4DEV:} & Gestión de riesgos \\
\textbf{Insumos:} & Acta de Constitución, cronograma, presupuesto y registro de interesados \\
\bottomrule
\end{tabularx}

\section{Enfoque}
Un riesgo es un evento incierto que, si ocurre, afecta el alcance, el tiempo, el costo o la calidad del proyecto. La gestión sigue cuatro pasos de PM4DEV:
\begin{enumerate}[leftmargin=*]
    \item \textbf{Identificar:} lluvia de ideas del equipo a partir de la EDT, el cronograma y el presupuesto. Cada riesgo se redacta como causa, evento y efecto.
    \item \textbf{Analizar:} se califican la probabilidad y el impacto con la escala de 1 a 3; la severidad es su producto.
    \item \textbf{Planificar la respuesta:} cada riesgo tiene una estrategia, una respuesta concreta, un disparador que avisa que está por ocurrir y un responsable.
    \item \textbf{Monitorear:} el registro se revisa cada martes en la reunión de control.
\end{enumerate}

\section{Escalas de Calificación}
{\small
\noindent\begin{tabularx}{\textwidth}{@{}C{1.1cm} P{2.9cm} L@{}}
\toprule
\rowcolor{AzulClaro}
\textbf{Valor} & \textbf{Probabilidad} & \textbf{Impacto} \\
\midrule
1 & Baja: hasta 20\,\% & Bajo: retraso menor a 3 días o costo menor a S/ 100; no afecta a los clientes. \\
2 & Media: de 21\,\% a 50\,\% & Medio: retraso de 3 a 7 días o costo de S/ 100 a S/ 300; afecta una función secundaria. \\
3 & Alta: más de 50\,\% & Alto: retraso mayor a una semana o costo mayor a S/ 300; pone en riesgo un hito, el piloto o un objetivo del Acta. \\
\bottomrule
\end{tabularx}
}

\vspace{0.15cm}
{\small
\noindent\begin{tabularx}{\textwidth}{@{}P{3.2cm} C{2.0cm} L@{}}
\toprule
\rowcolor{AzulClaro}
\textbf{Severidad (P $\times$ I)} & \textbf{Nivel} & \textbf{Tratamiento} \\
\midrule
6 a 9 & \etiqueta{SevAlta}{ALTA} & Respuesta activa desde ahora y revisión cada semana. \\
3 a 4 & \etiqueta{SevMedia}{MEDIA} & Respuesta planificada; se revisa cada semana su disparador. \\
1 a 2 & \etiqueta{SevBaja}{BAJA} & Se acepta y se revisa al cierre de cada fase. \\
\bottomrule
\end{tabularx}
}

\Needspace{14\baselineskip}
\section{Matriz de Severidad}
\begin{center}
\begin{tikzpicture}[x=2.6cm,y=1.25cm,font=\footnotesize]
\fill[SevBaja!22] (0,0) rectangle (1,1);
\draw[white, line width=1.5pt] (0,0) rectangle (1,1);
\node[anchor=north east, font=\tiny\bfseries, text=SevBaja] at (1-0.04,1-0.06) {1};
\fill[SevBaja!22] (1,0) rectangle (2,1);
\draw[white, line width=1.5pt] (1,0) rectangle (2,1);
\node[anchor=north east, font=\tiny\bfseries, text=SevBaja] at (2-0.04,1-0.06) {2};
\fill[SevMedia!22] (2,0) rectangle (3,1);
\draw[white, line width=1.5pt] (2,0) rectangle (3,1);
\node[anchor=north east, font=\tiny\bfseries, text=SevMedia] at (3-0.04,1-0.06) {3};
\node[font=\scriptsize\bfseries, align=center, text width=2.2cm] at (2.5,0.5) {R03};
\fill[SevBaja!22] (0,1) rectangle (1,2);
\draw[white, line width=1.5pt] (0,1) rectangle (1,2);
\node[anchor=north east, font=\tiny\bfseries, text=SevBaja] at (1-0.04,2-0.06) {2};
\node[font=\scriptsize\bfseries, align=center, text width=2.2cm] at (0.5,1.5) {R04};
\fill[SevMedia!22] (1,1) rectangle (2,2);
\draw[white, line width=1.5pt] (1,1) rectangle (2,2);
\node[anchor=north east, font=\tiny\bfseries, text=SevMedia] at (2-0.04,2-0.06) {4};
\node[font=\scriptsize\bfseries, align=center, text width=2.2cm] at (1.5,1.5) {R05 \quad R06};
\fill[SevAlta!30] (2,1) rectangle (3,2);
\draw[white, line width=1.5pt] (2,1) rectangle (3,2);
\node[anchor=north east, font=\tiny\bfseries, text=SevAlta] at (3-0.04,2-0.06) {6};
\node[font=\scriptsize\bfseries, align=center, text width=2.2cm] at (2.5,1.5) {R01 \quad R02 \quad R07};
\fill[SevMedia!22] (0,2) rectangle (1,3);
\draw[white, line width=1.5pt] (0,2) rectangle (1,3);
\node[anchor=north east, font=\tiny\bfseries, text=SevMedia] at (1-0.04,3-0.06) {3};
\fill[SevAlta!30] (1,2) rectangle (2,3);
\draw[white, line width=1.5pt] (1,2) rectangle (2,3);
\node[anchor=north east, font=\tiny\bfseries, text=SevAlta] at (2-0.04,3-0.06) {6};
\fill[SevAlta!30] (2,2) rectangle (3,3);
\draw[white, line width=1.5pt] (2,2) rectangle (3,3);
\node[anchor=north east, font=\tiny\bfseries, text=SevAlta] at (3-0.04,3-0.06) {9};
\node[below, font=\scriptsize] at (0.5,0) {Bajo (1)};
\node[left, font=\scriptsize] at (0,0.5) {Baja (1)};
\node[below, font=\scriptsize] at (1.5,0) {Medio (2)};
\node[left, font=\scriptsize] at (0,1.5) {Media (2)};
\node[below, font=\scriptsize] at (2.5,0) {Alto (3)};
\node[left, font=\scriptsize] at (0,2.5) {Alta (3)};
\node[below=16pt, font=\small\bfseries, text=AzulUNA] at (1.5,0) {Impacto};
\node[rotate=90, font=\small\bfseries, text=AzulUNA] at (-0.78,1.5) {Probabilidad};
\end{tikzpicture}
\end{center}

\section{Registro de Riesgos}
Ordenado de mayor a menor severidad. P: probabilidad; I: impacto; Sev.: severidad.

\vspace{0.1cm}
{\footnotesize
\setlength{\tabcolsep}{3pt}
\begin{longtable}{@{}P{0.75cm} >{\raggedright\arraybackslash}p{3.35cm} C{0.45cm} C{0.45cm} C{0.85cm} P{1.45cm} >{\raggedright\arraybackslash}p{3.9cm} >{\raggedright\arraybackslash}p{2.45cm}@{}}
\toprule
\rowcolor{AzulClaro}
\textbf{ID} & \textbf{Riesgo} & \textbf{P} & \textbf{I} & \textbf{Sev.} & \textbf{Estrategia} & \textbf{Respuesta} & \textbf{Disparador} \\
\midrule
\endfirsthead
\toprule
\rowcolor{AzulClaro}
\textbf{ID} & \textbf{Riesgo} & \textbf{P} & \textbf{I} & \textbf{Sev.} & \textbf{Estrategia} & \textbf{Respuesta} & \textbf{Disparador} \\
\midrule
\endhead
\bottomrule
\endfoot
\textbf{R01} & Si pocos propietarios aceptan publicar, el ranking tendrá pocos cuartos y el piloto perderá valor. & 2 & 3 & \etiqueta{SevAlta}{\,6\,} & Mitigar & Visitas de campo con la guía de publicación y acompañamiento por WhatsApp; dos visitas adicionales si es necesario. & Menos de 8 cuartos registrados al cierre de la semana 6. \\
\rowcolor{GrisFila}
\textbf{R02} & Si la migración de SQLite a Supabase toma más de lo previsto, se retrasan el registro, la nube y el MVP final. & 2 & 3 & \etiqueta{SevAlta}{\,6\,} & Mitigar & Práctica guiada con Supabase antes de la semana 9 y prueba de conexión en un proyecto aparte; dos jornadas nocturnas si el retraso persiste. & Inicio de sesión sin funcionar al cierre de la semana 10. \\
\textbf{R07} & Si pocos arrendatarios se inscriben en el piloto, no se podrá medir la satisfacción con la muestra prevista. & 2 & 3 & \etiqueta{SevAlta}{\,6\,} & Mitigar & Convocatoria en grupos de WhatsApp de estudiantes, afiches en la ciudad universitaria e incentivos para los participantes. & Menos de 10 arrendatarios inscritos al cierre de la semana 12. \\
\rowcolor{GrisFila}
\textbf{R05} & Si un integrante enferma o se ausenta, sus tareas se retrasan. & 2 & 2 & \etiqueta{SevMedia}{\,4\,} & Mitigar & Cada tarea queda documentada en su Issue y tiene un integrante de apoyo que puede continuarla. & Ausencia mayor a una semana. \\
\textbf{R06} & Si Supabase o el servicio de mapas abiertos limitan su plan gratuito, la app puede dejar de funcionar o generar costos. & 2 & 2 & \etiqueta{SevMedia}{\,4\,} & Mitigar & Fotos comprimidas antes de subirlas, uso moderado de los mapas, actividad semanal en la base de datos y revisión mensual de los términos de uso. & Uso mayor al 70\,\% de algún límite del plan gratuito. \\
\rowcolor{GrisFila}
\textbf{R03} & Si una laptop o el celular de pruebas falla o es robado, se pierden días de trabajo. & 1 & 3 & \etiqueta{SevMedia}{\,3\,} & Mitigar & Todo el código y los documentos se suben a GitHub cada día; se repara el equipo o se alquila uno de reemplazo. & Equipo inutilizado por más de 2 días. \\
\textbf{R04} & Si suben los precios de alimentación y transporte, el gasto en efectivo supera el aporte previsto. & 2 & 1 & \etiqueta{SevBaja}{\,2\,} & Aceptar & Se cubre con la contingencia y se compara cada mes el gasto real con el previsto. & Gasto mensual 10\,\% mayor que el previsto. \\
\end{longtable}
}

\Needspace{14\baselineskip}
\section{Responsables y Presupuesto de Contingencia}
La contingencia de cada riesgo es su valor esperado (probabilidad $\times$ impacto en soles), calculado en el presupuesto del proyecto. Los riesgos sin costo se tratan con acciones del equipo; si llegan a requerir dinero, se usa la reserva de gestión (S/ 1,166.35) con una solicitud de cambio.

\vspace{0.1cm}
{\small
\rowcolors{2}{white}{GrisFila}
\noindent\begin{tabularx}{\textwidth}{@{}C{0.8cm} P{1.7cm} L Q{2.2cm}@{}}
\toprule
\rowcolor{AzulClaro}
\textbf{ID} & \textbf{Categoría} & \textbf{Responsable del riesgo} & \textbf{Contingencia S/} \\
\midrule
R01 & Clientes & Clyde N. Paricahua Pari & 150.00 \\
R02 & Técnico & Eliseo Tarqui Ajahuana & 200.72 \\
R03 & Recursos & Viviana M. Quisbert Quispe & 80.00 \\
R04 & Costos & Ruth K. Apaza Solis & 68.60 \\
R05 & Equipo & Clyde N. Paricahua Pari & 56.08 \\
R06 & Técnico & Eliseo Tarqui Ajahuana & \textcolor{GrisOscuro}{sin costo} \\
R07 & Clientes & Viviana M. Quisbert Quispe & \textcolor{GrisOscuro}{sin costo} \\
\midrule
\rowcolor{white}\multicolumn{3}{@{}l}{\textbf{Reserva de contingencia}} & \textbf{555.40} \\
\bottomrule
\end{tabularx}
}

\section{Estrategias de Respuesta}
{\small
\rowcolors{2}{white}{GrisFila}
\noindent\begin{tabularx}{\textwidth}{@{}P{2.2cm} L@{}}
\toprule
\rowcolor{AzulClaro}
\textbf{Estrategia} & \textbf{Cuándo se usa} \\
\midrule
Evitar & Se cambia el plan para eliminar la causa del riesgo. \\
Mitigar & Se actúa antes para bajar la probabilidad o el impacto. \\
Transferir & Se traslada el impacto a un tercero, por ejemplo con una garantía. \\
Aceptar & Se asume el riesgo y se cubre con la contingencia si ocurre. \\
\bottomrule
\end{tabularx}
}

\section{Monitoreo y Control}
\begin{itemize}[leftmargin=*]
    \item \textbf{Martes:} la Gestora revisa los disparadores; si uno se activa, el responsable ejecuta la respuesta y se registra en el acta.
    \item \textbf{Nuevos riesgos:} cualquier integrante puede proponerlos en un Issue de GitHub con la etiqueta \texttt{riesgo}; se califican en la siguiente reunión.
    \item \textbf{Uso de la contingencia:} se anota el riesgo, el monto y el comprobante; la Analista de Planificación actualiza el saldo.
    \item \textbf{Cierre de un riesgo:} cuando su causa desaparece o pasa su ventana de tiempo, se marca como cerrado.
    \item \textbf{Actualización:} el registro se versiona en el repositorio al cierre de cada fase.
\end{itemize}

\Needspace{24\baselineskip}
\section{Firmas}
\noindent En señal de conformidad con el presente documento, los integrantes del equipo firman a continuación:

\vspace{1.5cm}
\begin{center}
\begin{tabular}{c c c}
    \parbox{6.5cm}{\centering \rule{6.0cm}{0.6pt}\\[5pt] \textbf{Clyde Neil Paricahua Pari}\\[2pt] {\footnotesize Project Manager}}
    & \hspace{0.8cm} &
    \parbox{6.5cm}{\centering \rule{6.0cm}{0.6pt}\\[5pt] \textbf{Ruth Karina Apaza Solis}\\[2pt] {\footnotesize Analista de Planificación}}
    \\[2.0cm]
    \parbox{6.5cm}{\centering \rule{6.0cm}{0.6pt}\\[5pt] \textbf{Eliseo Tarqui Ajahuana}\\[2pt] {\footnotesize Especialista Técnico / Arquitectura}}
    & \hspace{0.8cm} &
    \parbox{6.5cm}{\centering \rule{6.0cm}{0.6pt}\\[5pt] \textbf{Viviana Mary Quisbert Quispe}\\[2pt] {\footnotesize Gestora de Comunicación, Riesgos y Documentación}}
\end{tabular}
\end{center}
\end{document}
